% Converted from the authoritative recovery source, 2026-11-14.
\subsection{An atom-independent influence theorem for product
blocks}\label{influence-recovery:an-atom-independent-influence-theorem-for-product-blocks}

Let \(\nu\) be an arbitrary product law on Boolean blocks \(Z_a\), each
of size at most \(k\). Zeros and arbitrarily small atom probabilities
are permitted. A function has \textbf{block degree at most \(\Delta\)}
if it is a sum of functions each depending on at most \(\Delta\) blocks.
This is a pointwise representation on the entire Boolean cube, including
points outside the support of \(\nu\).

For \(g:\{-1,+1\}^N\to[0,1]\) of block degree at most \(\Delta\ge1\),
define

\[
 J_a^\nu(g)=\mathbb E_\nu\operatorname{Var}_{Z_a}
                      (g\mid Z_{-a})
 =\frac12\mathbb E(g(Z)-g(Z^{a\leftarrow Z_a'}))^2,
 \tag{B1}
\]

where \(Z_a'\) is an independent block with its original law. Write
\(v=\operatorname{Var}_\nu(g)>0\) and \(L=k+\log(\Delta/v)\), using
natural logarithms.

\textbf{Theorem.} An absolute \(c>0\) satisfies

\[
 \boxed{\quad
 \max_a J_a^\nu(g)
 \ge c\,\frac{v^{5/2}}
 {\Delta^6\,[k+\log(\Delta/v)]^6}.
 \quad} \tag{B2}
\]

This is independent of the number of blocks and their smallest atom
masses. Its influence resamples an entire block, rather than physically
flipping one bit in a possibly constrained block.

\textbf{Proof.} For \(\eta\in(0,1)\), round each block law \(\nu_a\) to
a dyadic law on \(b=k+\lceil\log_2(1/\eta)\rceil\) fair source bits.
Allocate at least \(\lceil(1-\eta)2^b\nu_a(z)\rceil\) source words to
each supported \(z\). These counts sum to at most \(2^b\); fill the
remaining words arbitrarily. The resulting law is

\[
 \nu'_a=(1-\eta)\nu_a+\eta r_a.
 \tag{B3}
\]

We compare moments by a single univariate interpolation. Put
\(\nu_a(t)=(1-t)\nu_a+t r_a\) and \(\nu(t)=\bigotimes_a\nu_a(t)\), for
\(0\le t\le1\). For any bounded function \(H\) of block degree at most
\(m\), \(\mathbb E_{\nu(t)}H\) is a polynomial in \(t\) of degree at
most \(m\). The univariate Markov inequality on \([0,1]\) bounds its
derivative by \(2m^2\) times its supremum absolute value. Applied to
\(g\) and \(g^2\), whose block degrees are at most \(\Delta\) and
\(2\Delta\), this gives

\[
 |\operatorname{Var}_{\nu'}g-\operatorname{Var}_{\nu}g|
 \le12\Delta^2\eta.
 \tag{B4}
\]

Indeed, the first moment changes by at most \(2\Delta^2\eta\), the
second by at most \(8\Delta^2\eta\), and squaring a \([0,1]\)-valued
first moment costs at most twice its change. Thus
\(\operatorname{Var}_{\nu'}(g)\ge v/4\) for sufficiently small
\(\Delta^2\eta/v\). Compose \(g\) with the independent dyadic block maps
to get \(h\) on fair bits. Each term depending on at most \(\Delta\)
blocks now depends on at most \(b\Delta\) fair bits. Therefore
\(\deg(h)\le b\Delta\) and \(h\) is globally bounded. By (U), some fair
bit \(j\) has

\[
 I_j(h)\ge \frac{c_0}{32}
             \frac{v^{5/2}}{b^6\Delta^6}.
 \tag{B5}
\]

Suppose \(j\) lies in the fair source of block \(a\). Conditional on all
other blocks, flipping \(j\) gives two identically distributed
\(g\)-values, with common variance \(\operatorname{Var}_{\nu'_a}(g)\).
For any identically distributed \(A,B\),
\(\mathbb E(A-B)^2\le4\operatorname{Var}(A)\). Hence

\[
 I_j(h)\le J_a^{\nu'}(g).
 \tag{B6}
\]

We next compare block influences without comparing the global densities.
Put \(H_a(z_{-a})=\operatorname{Var}_{\nu_a}(g(z_{-a},Z_a))\). This is a
globally bounded function in \([0,1/4]\) of block degree at most
\(2\Delta\). Replacing \(\nu_a\) by the mixture (B3) increases this
conditional variance by at most \(2\eta\): the mixture-variance identity
gives the sharper upper bound \(5\eta/4\). Apply the same univariate
interpolation to \(H_a\) while varying the other block laws. Markov
bounds the change in its expectation by \(2\Delta^2\eta\). Consequently

\[
 J_a^{\nu'}(g)
 \le J_a^\nu(g)+4\Delta^2\eta
 \tag{B7}
\]

This step remains valid with zero atoms and does not use any comparison
between the global product densities.

Choose

\[
 \eta=\alpha\frac{v^3}{\Delta^8 L^6},
 \qquad L=k+\log(\Delta/v),
 \tag{B8}
\]

where \(\alpha>0\) is an absolute sufficiently small constant. Then
\(b\le C_\alpha L\) with \(C_\alpha=O(1+\log(1/\alpha))\). The variance
condition in (B4) holds for small enough \(\alpha\). The error in (B7)
is at most \(4\alpha v^3/(\Delta^6L^6)\). The lower bound (B5) is at
least \(c_0v^{5/2}/(32C_\alpha^6\Delta^6L^6)\). Choose \(\alpha\) so
small that \(4\alpha\le c_0/(64C_\alpha^6)\), which is possible because
\(\alpha(1+\log(1/\alpha))^6\to0\). Since \(v\le1\), subtracting (B7)
proves (B2). This finishes the proof.
